% ============================================================
\section{Conclusión: el Teorema 4.2 completo}
% ============================================================

\begin{frame}{Teorema 4.2 — Resumen de la demostración}
  \begin{block}{Teorema 4.2}
    $\pione(X,\xo)$ es un \textbf{grupo} bajo la operación
    $\clase{\alpha}\circ\clase{\beta}=\clase{\alpha*\beta}$.
  \end{block}

  \bigskip
  \begin{columns}[T]
    \begin{column}{0.30\textwidth}
      \centering
      \textcolor{verdeTeal}{\textbf{Lema A}}\\[4pt]
      \textit{Identidad}\\[6pt]
      $\clase{c}\circ\clase{\alpha}=\clase{\alpha}$\\[4pt]
      Homotopía: zona constante se contrae hacia $t=0$
    \end{column}
    \begin{column}{0.33\textwidth}
      \centering
      \textcolor{morado}{\textbf{Lema B}}\\[4pt]
      \textit{Inversos}\\[6pt]
      $\clase{\alpha}\circ\clase{\abar}=\clase{c}$\\[4pt]
      Homotopía: avanza hasta $\alpha(s)$ y regresa
    \end{column}
    \begin{column}{0.33\textwidth}
      \centering
      \textcolor{coralOscuro}{\textbf{Lema C}}\\[4pt]
      \textit{Asociatividad}\\[6pt]
      $([\alpha]\circ[\beta])\circ[\gamma]$\\$=[\alpha]\circ([\beta]\circ[\gamma])$\\[4pt]
      Homotopía: desliza los puntos de corte
    \end{column}
  \end{columns}

  \bigskip
  \begin{center}
    \textbf{Técnica unificadora:} construir homotopías a partir de diagramas
    geométricos en el cuadrado $\IxI$, garantizando continuidad mediante el
    Lema de Continuidad.
  \end{center}
\end{frame}
